\section{Contrato del enlace y formato UART 8N1}

En el TP1, cada operando llegaba a la ALU mediante ocho señales paralelas: una
por cada bit. En este trabajo, la computadora envía esos mismos ocho bits de
forma secuencial mediante UART. Los bits de cada byte ingresan uno después de
otro por RX y la respuesta se transmite del mismo modo por TX. Como UART no
incorpora una señal de reloj, la computadora y la FPGA deben utilizar la misma
velocidad y el mismo formato de trama. En este diseño se emplean
\num{19200} baud
y el formato 8N1: ocho bits de datos, sin paridad y con un bit de stop.

Cada byte ocupa diez períodos de bit. Mientras no hay una transmisión, RX y TX
permanecen en nivel alto. Un bit de start en bajo marca el comienzo de la trama;
a continuación se envían los ocho bits de datos, desde el menos significativo
hasta el más significativo, y un bit de stop en alto indica el final.

\begin{figure}[H]
	\centering
	\begin{tikzpicture}[
			x=1.05cm,y=0.8cm,
			lbl/.style={font=\small,align=center}
		]
		\draw[->,thick] (-0.35,0) -- (10.55,0);
		\foreach \x in {0,...,10} {
			\draw[gray!55] (\x,-0.15) -- (\x,1.2);
		}
		\draw[very thick,codeviolet]
		(-0.3,1) -- (0,1) -- (0,0.2) -- (1,0.2);
		\draw[very thick,codeviolet]
		(1,0.2) -- (1,0.82) -- (2,0.82) -- (2,0.2) --
		(3,0.2) -- (3,0.82) -- (4,0.82) -- (4,0.2) --
		(5,0.2) -- (5,0.82) -- (6,0.82) -- (6,0.2) --
		(7,0.2) -- (7,0.82) -- (8,0.82) -- (8,0.2) --
		(9,0.2) -- (9,1) -- (10.4,1);
		\node[lbl] at (0.5,-0.42) {start};
		\foreach \x in {0,...,7} {
			\node[lbl] at ({1.5+\x},-0.42) {D\x};
		}
		\node[lbl] at (9.5,-0.42) {stop};
		\draw[<->,codeblue] (0,1.37) -- node[above,lbl]
		{$T_{\mathrm{bit}}$} (1,1.37);
	\end{tikzpicture}
	\caption{Estructura temporal de una trama UART 8N1.}%
	\label{fig:trama-uart}
\end{figure}

Las líneas verticales de la figura~\ref{fig:trama-uart} delimitan diez
intervalos iguales. El primer nivel bajo no es parte del dato: permite al
receptor detectar el inicio. D0 llega antes que D7, por lo que el orden temporal
no coincide con la forma habitual de escribir un byte. El stop alto confirma el
final y devuelve la línea a su nivel de reposo.

A \num{19200} baud, un bit dura aproximadamente
\SI{52.083}{\micro\second} y cada byte ocupa \SI{520.833}{\micro\second}.
El envío de los tres bytes de entrada demanda al menos
\SI{1.563}{\milli\second}; la respuesta agrega otro intervalo de byte. Estos
valores no incluyen demoras del sistema operativo o del adaptador USB.

\section{Generación del tick de muestreo}

El receptor utiliza sobremuestreo 16x: por cada bit recibido se generan
dieciséis pulsos de muestreo. De este modo, cada período de bit queda dividido
en dieciséis intervalos que sirven como referencia temporal. Después de
detectar la transición que puede indicar un bit de start, el receptor espera
ocho pulsos para comprobar su valor cerca del centro. Si el start es válido,
toma una muestra de cada bit de datos cada dieciséis pulsos y finalmente hace
lo mismo con el bit de stop.

Para obtener esos dieciséis pulsos por bit, el generador produce
\texttt{s\_tick} a una frecuencia 16 veces mayor que la tasa UART:

\begin{equation}
	f_{\mathrm{tick}} = 16 \cdot 19200 =
	\SI{307200}{\hertz}.
\end{equation}

Como el reloj de la placa es de \SI{100}{\mega\hertz}, se utiliza un divisor
entero redondeado al valor más próximo:

\begin{equation}
	D = \operatorname{round}
	\left(\frac{100\,000\,000}{307\,200}\right) = 326.
\end{equation}

El resultado no es exacto. El tick real es
\SI{306748.466}{\hertz} y la tasa efectiva es
\SI{19171.779}{baud}, con un error relativo de aproximadamente
\SI{-0.147}{\percent}. La diferencia acumulada a lo largo de una trama es muy
inferior a medio bit y queda dentro del margen que ofrece el muestreo central.

El RTL conserva el reloj de \SI{100}{\mega\hertz} y genera un pulso de
habilitación cada 326 ciclos. La suma de media frecuencia antes de la división
implementa el redondeo entero.

\begin{codebox}[language=SystemVerilog]
	localparam integer TICK_FREQ_HZ = BAUD_RATE * OVERSAMPLE;
	localparam integer DIVISOR =
	(CLK_FREQ_HZ + (TICK_FREQ_HZ / 2)) / TICK_FREQ_HZ;

	always @(posedge clk) begin
	if (reset) begin
	counter <= '0;
	s_tick  <= 1'b0;
	end else if (counter == DIVISOR - 1) begin
	counter <= '0;
	s_tick  <= 1'b1;
	end else begin
	counter <= counter + 1'b1;
	s_tick  <= 1'b0;
	end
	end
\end{codebox}

El fragmento muestra qué hardware se obtiene: un contador, un comparador contra
325 y un registro para el pulso. Como \texttt{s\_tick} permanece bajo en todos
los demás ciclos, RX y TX avanzan sus contadores de sub-bit únicamente cuando
corresponde, pero todos sus registros siguen capturados por \texttt{clk}.

\section{Protocolo de aplicación}

Cada operación tiene longitud fija. La computadora transmite A, luego B y por
último el opcode. Una vez recibidos los tres bytes válidos, la FPGA responde con
un byte de resultado.

\begin{table}[H]
	\centering
	\begin{tabularx}{0.94\textwidth}{ZZZZZ}
		\toprule
		\textbf{Sentido} & \textbf{Byte 0} & \textbf{Byte 1} &
		\textbf{Byte 2} & \textbf{Respuesta} \\
		\midrule
		PC a FPGA & A, 8 bits & B, 8 bits &
		Opcode en bits 5--0 & --- \\
		FPGA a PC & --- & --- & --- & Resultado, 8 bits \\
		\bottomrule
	\end{tabularx}
	\caption{Orden de los bytes del protocolo binario.}%
\end{table}

Los bits 7 y 6 del tercer byte se ignoran. Esta elección permite enviar el
opcode dentro de una palabra UART completa sin agregar codificación ASCII. El
protocolo es deliberadamente mínimo: no se devuelve un byte separado de estado
porque el requisito funcional solicita el resultado y las banderas continúan
visibles en LED8, LED9 y LED10.

\begin{table}[H]
	\centering
	\begin{tabular}{lccc}
		\toprule
		\textbf{Operación} & \textbf{Opcode} & \textbf{Hexadecimal} &
		\textbf{Resultado} \\
		\midrule
		SRA & \texttt{000011} & \texttt{03} & A desplazado aritméticamente
		B lugares \\
		SRL & \texttt{000010} & \texttt{02} & A desplazado lógicamente B lugares \\
		ADD & \texttt{100000} & \texttt{20} & A + B \\
		SUB & \texttt{100010} & \texttt{22} & A - B \\
		AND & \texttt{100100} & \texttt{24} & A AND B \\
		OR  & \texttt{100101} & \texttt{25} & A OR B \\
		XOR & \texttt{100110} & \texttt{26} & A XOR B \\
		NOR & \texttt{100111} & \texttt{27} & NOT (A OR B) \\
		\bottomrule
	\end{tabular}
	\caption{Operaciones conservadas del TP1.}%
\end{table}

Por ejemplo, para calcular \texttt{7F + 01}, el host transmite
\texttt{7F 01 20}. La respuesta esperada es \texttt{80}; al mismo tiempo,
LED10 se enciende para indicar overflow y los LED de carry y cero permanecen
apagados.

\section{Recuperación y casos inválidos}

El controlador debe volver a su estado inicial si la computadora deja
de enviar datos antes de completar una solicitud. Una vez recibido A,
o A y B, espera como máximo \SI{10}{\milli\second} el byte siguiente.
Si vence ese plazo, descarta la operación parcial y el próximo byte
se interpreta nuevamente como A. El timeout supera ampliamente los
\SI{0.521}{\milli\second} de una trama normal, aun cuando los bytes
se envían de manera individual.

\begin{table}[H]
	\centering
	\begin{tabularx}{\textwidth}{L{4.0cm}L{4.6cm}Y}
		\toprule
		\textbf{Situación} & \textbf{Acción} & \textbf{Respuesta} \\
		\midrule
		Start vuelve a alto antes del centro & Abortar la trama y esperar
		reposo alto. & Ninguna. \\
		Stop bajo & Emitir error de trama y esperar reposo alto. & Ninguna. \\
		FIFO RX lleno & Conservar los cuatro bytes previos y notificar
		overrun. & Se aborta la transacción parcial. \\
		Transacción incompleta & Vencer el contador a los
		\SI{10}{\milli\second}. & Ninguna. \\
		Opcode no reconocido & Ejecutar el caso por defecto de la ALU. &
		Un byte \texttt{00}, con Z activo. \\
		FIFO TX lleno & Conservar resultado e indicadores. & Esperar espacio
		antes de encolar. \\
		\bottomrule
	\end{tabularx}
	\caption{Política de tratamiento de condiciones excepcionales.}%
\end{table}
